\section{Epistemologi Nilai Kebenaran: Fakta Empiris dan Derivasi Matematis}

Setelah suatu ujaran ditetapkan sebagai \textbf{proposisi}, langkah berikutnya adalah menentukan nilai kebenarannya: \textbf{benar (T)} atau \textbf{salah (F)}. Nilai tersebut sering dinotasikan dengan fungsi nilai kebenaran; misalnya untuk simbol proposisi $p$, penulisan $\tau(p) = \mathbf{T}$ berarti ``proposisi $p$ bernilai benar'' \cite{rosen2019}.

Cara menetapkan apakah suatu proposisi benar atau salah dapat dibedakan secara epistemologis, terutama antara pengetahuan yang bersandar pada \textbf{observasi empiris} dan pengetahuan yang bersandar pada \textbf{derivasi deduktif} (rasionalisme matematis).

\subsection{Dasar Observasi Empiris}

Banyak validasi dalam kehidupan sehari-hari---termasuk pengujian hipotesis ilmiah---menggunakan pengamatan terhadap fakta di dunia nyata. Nilai kebenaran T/F tidak ``diturunkan'' dari intuisi semata, melainkan dapat didukung oleh pengukuran, eksperimen, atau dokumentasi observasi.

\begin{contoh}[Validasi empiris]
Diberikan proposisi \\
\indent $q$ : \textit{``Air pada tekanan atmosfer standar di wilayah pegunungan Himalaya tidak mendidih pada suhu $80^\circ$ Celsius.''}\\
Untuk menetapkan $\tau(q)$, digunakan pengamatan dan pengukuran (misalnya percobaan pemanasan air) pada kondisi yang relevan, bukan sekadar kesimpulan dari rumus tanpa konteks.
\end{contoh}

\subsection{Dasar Derivasi Tak-Empiris (Rasionalisme Matematis)}

Di bidang matematika dan sejumlah cabang ilmu komputer, kebenaran banyak proposisi ditetapkan melalui \textbf{deduksi dari definisi, aksioma, dan aturan inferensi}, tanpa melibatkan pengukuran fisik langsung \cite{wiki_first_order_logic}.

\begin{contoh}[Validasi melalui derivasi aljabar]
Perhatikan proposisi \\
\indent $r$ : \textit{``Penyelesaian persamaan $3x - 1 = 5$ adalah $x = 2$.''}\\
Nilai kebenaran $\tau(r) = \mathbf{T}$ ditunjukkan secara deduktif:
\begin{align*}
    3x - 1 &= 5 \\
    3x &= 6 \\
    x &= 2
\end{align*}
Tanpa perlu observasi lapangan, kesimpulan dapat diperoleh secara konsisten dari manipulasi aljabar yang sah.
\end{contoh}

Membedakan jalur empiris dan jalur deduktif membantu dalam merancang verifikasi: data di lapangan untuk hipotesis observasional, dan pembuktian formal untuk pernyataan matematis atau spesifikasi logika yang mengikuti aturan tertutup.
